\chapter{Implementation Log and Lessons Learned}
\renewcommand{\thefigure}{C.\arabic{figure}}
\renewcommand{\thetable}{C.\arabic{table}}
\label{app:lessons}

This appendix records the changes made to the training system after the thesis experiments,
each with the evidence that motivated it and what it taught. Entries are kept in the order the
work was done, and every entry follows the same pattern: the change, why it was made, and the
lesson. New entries are added at the end.

% ------------------------------------------------------------------------------------------
\section{Narrowing the C51 support to $\pm 2$}

\paragraph{Change.} The value support of the distributional head \cite{C51} was narrowed from
$[-10, 10]$ to $[-2, 2]$ (51 atoms, so 0.08 per atom instead of 0.4).

\paragraph{Why.} A return is the outcome ($+1$, $-1$ or the draw reward $-0.5$) plus the
potential-based shaping $\gamma\Phi(s') - \Phi(s)$ \cite{NgPBRS}, which telescopes over a game to
$-\Phi(s_0)$. With $\Phi \in [-0.4, 1]$ every return therefore lies in $[-2, 1.4]$; the returns
observed in the formal\_2 run spanned $-1.54$ to $0.71$. The $\pm 10$ support was inherited from
the thesis configuration, whose shaping was not potential-based. On it, the shaping of a capture
(about 0.02 to 0.05) was a tenth of an atom: the head could hardly represent the difference
between a good and a bad exchange.

\paragraph{Lesson.} The support of a distributional critic is part of the reward design, not a
default. It should be derived from the bounds of the return (here the telescoping of the
shaping makes that exact) and checked against the returns a run actually produces.

% ------------------------------------------------------------------------------------------
\section{Measuring the average Q-value on the actions taken}

\paragraph{Change.} The logged ``average Q'' is now the expected value of the action that was
taken in each sampled replay transition.

\paragraph{Why.} The thesis metric took the maximum over all 3600 actions. Nearly all of them are
illegal in any given state and are never trained, so the maximum was an artefact: it read about
$+2$ while the returns were about $-1$.

\paragraph{Lesson.} With a large, mostly illegal action space, any statistic taken over all
actions measures the untrained part of the network. Diagnostics must be restricted to actions
that can occur.

% ------------------------------------------------------------------------------------------
\section{Diagnosing a flat value function}

\paragraph{Observation.} After thousands of episodes of the formal\_2 run the agent still drew
most games, and its Q-values barely distinguished good moves from bad ones.

\paragraph{Diagnosis.} Two probes separated the possible causes. A \emph{clean-signal probe}
trained the same network on synthetic transitions in which only the combat outcome determined the
reward; it learned the difference quickly, so the learner itself (loss, target, optimiser) was
sound. A \emph{same-state probe} then compared, in real positions, the values of a winning and a
losing attack available in the same state. In the formal\_3 run at episode 1250 the winning attack
was preferred in 68\% of such positions, but by a margin of 0.003, against a true difference of
about 0.05. In real training the combat signal is small next to the noise: a winning attack is
worth $+0.019$ of shaping and a losing one $-0.033$, while an ordinary move's shaping varies with
a standard deviation of 0.02, and the outcome arrives hundreds of moves later.

\paragraph{Lesson.} When a reinforcement learner does not learn, first test the learner on a
problem whose answer is known; if it passes, the problem is in the signal (reward scale, support,
credit assignment), not in the optimisation. Probes that compare actions \emph{within the same
state} are far more informative than averages over states.

\paragraph{Operational note.} A run that ends with exit code 1 and no Python traceback was
killed from outside the process (here by memory pressure), not by a bug in the training code.

% ------------------------------------------------------------------------------------------
\section{Making the training loop faster}

\paragraph{Changes.}
\begin{itemize}
  \item \textbf{Batched evaluation.} The evaluation games now run side by side, with one network
        call per side per round instead of one call per position. An evaluation of
        $40 + 40$ games went from 363\,s to 143\,s.
  \item \textbf{Cheaper replay.} The C51 loss needs the distribution of the action taken only,
        and the target needs that of the Double-DQN action only. The network now returns raw
        logits; the two needed actions are gathered first and normalised afterwards, instead of
        normalising all $3600 \times 51$ outputs. One update went from 128.7\,ms to 97.5\,ms.
  \item \textbf{More lanes at a constant replay ratio.} The loop used to perform at most one
        gradient update per iteration, so adding parallel games silently lowered the number of
        updates per move. It now performs one update per \texttt{REPLAY\_UPDATE\_INTERVAL}
        learner moves however many lanes moved, which allowed 16 lanes instead of 8 (acting cost
        per move fell from 2.0\,ms to 1.33\,ms).
\end{itemize}

\paragraph{Lesson.} Profile before optimising: replay was half of the step time and evaluation
was serial, while the environment was below 10\%. And a ``speed'' setting such as the number of
lanes can change the learning algorithm (here the replay ratio) unless the loop is written so
that it cannot.

% ------------------------------------------------------------------------------------------
\section{Removing the curriculum's environment changes}

\paragraph{Change.} The five-phase curriculum no longer changes the game. Every stage is played
under fog of war, with the same turn limit (1500), the same no-battle draw rule (200 plies,
also used in evaluation and the GUI) and the same, constant shaping weight. The stages remain
only as an opponent schedule: random and scripted opponents first, then the stronger heuristic,
then self-play and the league, advancing on evaluation gates as before.

\paragraph{Why.} The thesis curriculum \cite{CurriculumLearning} changed the environment under
the network at every transition. Phase 1 had full observability (the agent saw every enemy
rank), and the shaping weight fell from 1.0 to 0.5, 0.2 and finally 0. Each of these changes the
optimal values of the states the network has already learned: a value learned with the enemy's
ranks visible is wrong under fog of war, and cutting the shaping weight changes every return.
The network must then relearn what it knew, with the replay memory still full of transitions
from the old rules, which is the setting in which catastrophic forgetting
\cite{CatastrophicForgetting} occurs. Annealing the shaping also buys nothing: potential-based
shaping leaves the optimal policy unchanged \cite{NgPBRS}, so there is no bias to remove. A
smaller inconsistency was the draw rule: training drew after 100 plies without a battle and
evaluation after 200, so the agent was graded under a rule it had not trained under.

\paragraph{Lesson.} A curriculum should make the \emph{task} easier (weaker opponents), not
change the \emph{rules} the values refer to. Anything the network's targets depend on
(observability, reward, termination) must stay fixed across stages, and evaluation must use the
same rules as training.

% ------------------------------------------------------------------------------------------
\section{Cold start from demonstrations}

\paragraph{Change.} Before its first game the agent now learns from games in which the scripted
heuristic plays its side, following Deep Q-learning from Demonstrations (DQfD) \cite{DQfD}; the
idea of a short supervised ``cold start'' before reinforcement learning is also central to recent
reasoning-model training \cite{DeepSeekR1}.
\begin{enumerate}
  \item The heuristic plays 100 games against the random and heuristic opponents under the
        training rules. Its transitions are recorded exactly as the agent's own are (move plus
        reply, PBRS shaping plus outcome, 5-step returns, mirror augmentation) into a separate
        demonstration buffer that the agent's own experience never overwrites. Each one also
        keeps the legal moves of its state.
  \item The network is pre-trained on the demonstrations alone (5000 updates).
  \item During training 10\% of each batch comes from the demonstrations, falling linearly to
        0\% by episode 5000.
  \item Demonstrated transitions also get the large-margin loss of DQfD,
        \[
          J_E(Q) = \max_{a \in \mathcal{A}(s)} \bigl[ Q(s,a) + \ell(a_E, a) \bigr] - Q(s, a_E),
          \qquad \ell(a_E, a) = \begin{cases} 0 & a = a_E \\ 0.1 & a \neq a_E \end{cases}
        \]
        with the maximum over the legal moves $\mathcal{A}(s)$ only, added to the C51 loss with
        weight 1.
\end{enumerate}
The earlier ``imitation'' option, which let the heuristic choose 10\% of the agent's moves in
its first 2000 games, is switched off.

\paragraph{Why.} A network trained from zero starts from random play, which against these
opponents mostly ends in draws: at episode 1000 the formal\_3 run drew 32 of 40 evaluation games
against the random opponent. Its first thousands of games therefore carry almost no outcome
signal. The imitation option did not fix this: the heuristic's moves entered the replay memory
like the agent's own, so the C51 loss learned what they were worth but nothing told the network
to prefer them over the moves that were not played. The margin loss does exactly that. In a unit
test with equal returns for every move, the C51 loss alone left the greedy move at chance, while
with the margin loss the demonstrated move became the greedy one in at least 90\% of states.
The margin is 0.1 (1.25 atoms of the $\pm 2$ support), comparable to the value of a capture; DQfD
used 0.8 with rewards clipped to $\pm 1$. The demonstration share decays because the heuristic is
a weak demonstrator, not an expert, and should not cap the agent.

\paragraph{Result.} The 100 demonstration games took about three minutes; the heuristic won
about half of them (52 wins, 26 losses, 22 draws in one collection) and they produced 38\,742
transitions, mirror copies included. Pre-training costs about 107\,ms per update, so the 5000
updates take about nine minutes. Whether the cold start speeds up learning has not yet been
measured in a full run.

% ------------------------------------------------------------------------------------------
\section{The formal\_3 run and the start of formal\_4}

\paragraph{formal\_3.} The formal\_3 run trained from zero with the $\pm 2$ support but still
with the curriculum's environment changes and without the cold start. It was stopped at episode
3500, never having left Phase~1. Its evaluations (40 greedy games per opponent every 500
episodes) did not improve; against the random opponent they got worse:
\begin{center}
\begin{tabular}{rcc}
\toprule
Episode & vs random (W/L/D) & vs heuristic (W/L/D) \\
\midrule
500  & 8/14/18 & 1/37/2  \\
1000 & 7/1/32  & 1/28/11 \\
1500 & 6/2/32  & 1/32/7  \\
2000 & 0/4/36  & 0/29/11 \\
2500 & 5/8/27  & 2/30/8  \\
3000 & 3/5/32  & 0/35/5  \\
3500 & 1/3/36  & 0/35/5  \\
\bottomrule
\end{tabular}
\end{center}
The agent learned to avoid losing to the random opponent mainly by drawing, and it never won
more than 2 of 40 games against the heuristic. Its games carried almost no outcome signal, which
is the problem the cold start addresses.

\paragraph{formal\_4.} formal\_4 starts from zero with every change above: the opponent schedule
without environment changes, the cold start, the asynchronous learner and 16 lanes. Its cold
start played the 100 demonstration games in 84\,s (the heuristic won 56, lost 24 and drew 20)
and stored 35\,544 transitions. The 5000 pre-training updates took 271\,s (54\,ms per update with
the GPU not shared with another run), and the large-margin loss fell to 0.068, below the margin
of 0.1.

\paragraph{formal\_4, first 1000 episodes.} The evaluations (40 greedy games per opponent) were,
next to formal\_3 at the same episodes:
\begin{center}
\begin{tabular}{rcccc}
\toprule
 & \multicolumn{2}{c}{formal\_4} & \multicolumn{2}{c}{formal\_3} \\
Episode & vs random & vs heuristic & vs random & vs heuristic \\
\midrule
500  & 16/16/8 & 10/17/13 & 8/14/18 & 1/37/2  \\
1000 & 14/22/4 & 8/21/11  & 7/1/32  & 1/28/11 \\
\bottomrule
\end{tabular}
\end{center}
(wins/losses/draws). Against the heuristic formal\_4 won 8 to 10 games where formal\_3 won one,
and it no longer drew most games against the random opponent. But it lost more than half of them
instead, almost always by losing all its movable pieces, and it did not improve from episode 500
to 1000, nor by episode 2500 (vs random 12/23/5, 12/25/3 and 15/22/3 at episodes 1500, 2000 and
2500; vs heuristic 4/14/22, 5/19/16 and 6/16/18). formal\_4 also ran about 1.4 times as fast (1000 episodes in 45 minutes, against about
64 for formal\_3).

\paragraph{Diagnosis: attacks on revealed pieces.} A probe played 30 greedy games against the
random opponent under the training rules and classified every attack by whether the attacked
piece's rank was already known to the attacker:
\begin{center}
\begin{tabular}{lccc}
\toprule
Player 1 & attacks per game & hidden target: won / lost & revealed target: won / lost \\
\midrule
heuristic (the demonstrator) & 39.1 & 39\% / 47\% & 77\% / 0\%  \\
formal\_4, episode 250       & 41.7 & 45\% / 45\% & 26\% / 62\% \\
formal\_4, episode 500       & 41.0 & 42\% / 43\% & 27\% / 63\% \\
formal\_4, episode 1000      & 39.8 & 43\% / 42\% & 28\% / 63\% \\
formal\_3, episode 1000      & 20.2 & 36\% / 48\% & 22\% / 72\% \\
\bottomrule
\end{tabular}
\end{center}
Against hidden pieces the agent does as well as the heuristic. Against pieces whose rank it can
see (the ranks of revealed enemy pieces are in the observation, and a test checks that they reach
the network) it loses most of its attacks, while the heuristic never attacks a known piece it
cannot beat. The cold start transferred how \emph{often} the heuristic attacks (about 40 attacks
per game, twice formal\_3's), but not \emph{which} known pieces it attacks, and this was so from
episode 250 on, so it was never learned rather than learned and forgotten.

\paragraph{Test: pre-training on the demonstrations alone.} To see whether more imitation would
teach the rule, a separate copy of the network was pre-trained on a fresh set of 100
demonstration games (35\,198 transitions) with the cold start's losses, for 20\,000 updates
instead of 5000, and probed in the same way. The agreement is the share of demonstration states
in which the network's greedy legal move is the heuristic's:
\begin{center}
\begin{tabular}{rcccc}
\toprule
Updates & agreement & margin loss & revealed target: won / lost & hidden target: won / lost \\
\midrule
0      & 8.4\%  & --    & --          & --          \\
5000   & 63.1\% & 0.072 & 28\% / 60\% & 41\% / 45\% \\
10\,000 & 75.7\% & 0.051 & 28\% / 60\% & 44\% / 42\% \\
20\,000 & 88.2\% & 0.027 & 31\% / 62\% & 42\% / 44\% \\
\bottomrule
\end{tabular}
\end{center}
The network came to reproduce the heuristic's move in 88\% of the demonstration states, yet in its
own games it lost as many attacks on revealed pieces as before. It also did worse against pieces
it could see than against hidden ones, as a player would that picks its targets without regard to
their visible rank: the enemy pieces that are revealed are mostly ones that survived an earlier
battle, so they are stronger than average. Fitting the demonstrations, even closely, did not
produce the rule that makes the heuristic's attacks good. (The agreement is measured on the
demonstrations the network was trained on, so part of it may be memorisation.)

\paragraph{Test: does the value of an attack respond to the target's rank?} In 300 real
positions (random play, fog of war) where a piece of rank Miner to General could attack an
adjacent enemy piece, the observation was edited at the target square only, showing the target
as a revealed Scout (the attack wins), a revealed Marshal (the attack loses) or as hidden, and the
network's value of that attack was compared across the three. The same edit was also applied
while reading the value of a quiet move far from the target, as a control.
\begin{center}
\begin{tabular}{lcccc}
\toprule
Network & $Q$ vs Scout & $Q$ vs Marshal & $Q$ vs hidden & Scout $-$ Marshal (positive in) \\
\midrule
untrained               & $-0.023$ & $-0.025$ & $-0.022$ & $+0.002$ (51\%) \\
formal\_3, episode 3500  & $-0.642$ & $-0.644$ & $-0.628$ & $+0.001$ (46\%) \\
formal\_4, episode 2500  & $+0.055$ & $+0.048$ & $+0.187$ & $+0.007$ (50\%) \\
pre-trained, 20\,000 updates & $-0.140$ & $-0.172$ & $+0.074$ & $+0.032$ (59\%) \\
\bottomrule
\end{tabular}
\end{center}
The control moved by less than 0.01 in every case. The networks trained by reinforcement value
an attack on a revealed Scout and on a revealed Marshal the same, no more apart than the untrained
network does, although one wins a piece and the other loses one. They do respond to \emph{whether}
the target is revealed (formal\_4 values attacking a hidden piece at $+0.19$ and a revealed one at
about $+0.05$), but not to \emph{which} rank is revealed. The pre-trained network shows a small
preference ($+0.03$), with a spread across positions (standard deviation 0.12) four times as
large. The environment keeps a revealed piece's rank visible when it moves, and the ranks are in
the observation, so the information is there: it is the training signal that has not made the
networks use it. This is the combat signal of the earlier diagnosis: the shaping difference
between winning and losing an exchange (a few hundredths) is no larger than the variation of
an ordinary move's shaping, and the outcome comes hundreds of moves later.

\paragraph{Measurement: the combat signal of the potential.} The shaping of Player-1
transitions (its move plus the reply, as the trainer builds them; terminal transitions left out)
was measured in 30 games against the random opponent, with the heuristic and with a random player
as Player 1, for the current potential and for its material term alone,
$\Phi_{\mathrm{mat}} = (\text{own} - \text{enemy material}) / 148$:
\begin{center}
\begin{tabular}{lcc}
\toprule
 & current $\Phi$ & material only \\
\midrule
quiet move, mean / sd (random Player 1)        & $+0.0008$ / 0.015 & $0.0000$ / 0.0001 \\
quiet move, mean (heuristic Player 1)          & $+0.0065$         & $0.0000$ \\
attack on a revealed piece, won $-$ lost (random) & 0.043          & 0.061 \\
\quad in standard deviations of a quiet move   & 2.9               & about 430 \\
\quad in atoms of the $\pm 2$ support          & 0.53              & 0.76 \\
attack on a hidden piece, won $-$ lost (heuristic) & 0.093         & 0.058 \\
\bottomrule
\end{tabular}
\end{center}
Under the current potential an ordinary move is not neutral: the proximity and penetration terms
reward advancing, so for the heuristic a quiet move is worth $+0.0065$ on average, and beating a
revealed piece only $+0.011$. Every 5-step return also adds four more transitions of this
variation (standard deviation about 0.02 each) to the difference between a won and a lost
exchange of about 0.04. With the material term alone, a move without a battle is worth exactly
nothing (up to the discount), so the only shaping left in a return is that of the battles in it,
and a won exchange differs from a lost one by 0.06, about three-quarters of an atom. The
information and proximity terms also explain why the networks prefer attacking hidden pieces:
revealing a piece raises the information term whatever the result.

\paragraph{Lesson.} A shaping potential with several heuristic terms makes every move carry
reward, and that variation is noise to the decision the agent most needs to learn. Before tuning
a learner, measure the reward difference between the decisions that matter, relative to the
variation of the rest of the reward stream, and in the units the critic resolves.

\paragraph{Lesson.} Imitation transfers the frequent, easy-to-see features of a demonstrator's
play (here, attacking a lot) before the conditional rule behind them (attack a known piece only
if it is weaker), and a high agreement with the demonstrator does not show that the rule was
learned. A demonstration is worth checking for the specific decisions it was meant to teach, with
a probe that classifies those decisions in the agent's own games, not only by win rate or by
agreement on the demonstrations.

\paragraph{Lesson.} An evaluation that only counts wins hides how a policy fails: formal\_3's
win rate against the random opponent stayed low throughout, but its draws rising from 18 to 36
of 40 showed that it was converging to passive play rather than failing to learn at all.

\paragraph{Lesson.} Demonstrations stored as ordinary experience teach values, not preferences;
to make a value learner imitate, the loss must rank the demonstrated action above the
alternatives. And with a large, mostly illegal action space, that ranking has to be restricted to
the legal moves, or it spends its effort pushing down actions that can never be played.

% ------------------------------------------------------------------------------------------
\section{A material-only potential (formal\_5)}

\paragraph{Change.} The weights of the shaping potential's four terms (material, penetration,
information, proximity) became a setting, \texttt{POTENTIAL\_WEIGHTS}, and were changed from
$(0.4, 0.15, 0.1, 0.35)$ to $(1, 0, 0, 0)$:
\[
\Phi(s) = \frac{\text{own material} - \text{enemy material}}{148},
\]
with a hidden enemy piece still counted at the mean army rank of 3.7. Its range is $[-1, 1]$, and
the support test now reads the range from the weights. $\Phi$ is 0 at the start of a game (both
armies have 148 material, the hidden one counted at the mean), so the shaping of a whole game
telescopes to zero, and a return still lies within the outcome reward plus or minus 1, inside the
$\pm 2$ support. Everything else is formal\_4's configuration, including the cold start and the
asynchronous learner. formal\_4 was stopped at episode 3000; its last evaluation had 17 wins,
17 losses and 6 draws against the random opponent and 3 wins, 24 losses and 13 draws against the
heuristic, with no gain since episode 1500.

\paragraph{Why.} The measurement above: under the old potential an ordinary move carried a
shaping reward that varied as much as the difference between winning and losing an exchange on a
revealed piece, and the network did not learn to tell those exchanges apart. With material only,
a move without a battle carries no shaping, and a won exchange is worth about three-quarters of
an atom more than a lost one. The cost is that nothing rewards advancing or scouting any more
until a battle happens; whether the agent still finds the enemy, or stalls into no-battle draws,
is what formal\_5 has to show.

\paragraph{Result, first 1000 episodes.} The cold start went as in formal\_4 (the heuristic won 58,
lost 20 and drew 22 of the 100 demonstration games; the margin loss fell to 0.071). The
evaluations were 20/16/4 and 9/30/1 against the random opponent and 3/18/19 and 9/19/12 against
the heuristic at episodes 500 and 1000 (wins/losses/draws), against 16/16/8 and 14/22/4, 10/17/13
and 8/21/11 for formal\_4. The agent did not stall: no-battle draws were 74 and 94 of the
training games in the two blocks of 500, and games lasted 340 to 346 moves. It still lost
mostly by running out of movable pieces (197 and 198 of 500 games), as formal\_4 did. So far the
material-only potential has not changed how the agent loses. Two evaluations 500 episodes apart
also differed by 11 wins out of 40 against the same opponent, which shows how wide the noise of a
40-game evaluation is.

\paragraph{Result, to episode 2750.} formal\_5 was stopped at episode 2750 to start formal\_6
(Section~\ref{sec:formal6}) with the fixes found in the meantime. Its later evaluations were
22/15/3, 17/22/1 and 16/20/4 against the random opponent and 9/24/7, 8/20/12 and 7/14/19 against
the heuristic at episodes 1500, 2000 and 2500. formal\_4 had 12, 12 and 15 wins against the random
opponent and 4, 5 and 6 against the heuristic at the same episodes. The material-only potential
won a few more games against the heuristic, but by less than the noise of 40-game evaluations
(about 3 wins), and neither run passed a gate. The potential was not the bottleneck it was
suspected to be; the faults of the next sections were found while formal\_5 ran.

\paragraph{Result, the revealed-piece probes.} The two probes of formal\_4 were repeated on the
networks of formal\_5 (episode 2750) and formal\_4 (episode 3000). Both have 30 input planes, so
they were run under the formal\_6 code with the agent built for 30 planes and given the first 30
planes of the new encoding, whose layout is unchanged. Each played 30 greedy games as Player~1
against the random opponent:
\begin{center}
\begin{tabular}{lccc}
\toprule
 & formal\_5, ep.\ 2750 & formal\_4, ep.\ 3000 & heuristic \\
\midrule
games won / lost / drawn               & 17 / 13 / 0 & 14 / 16 / 0 & 23 / 4 / 3 \\
attacks on revealed pieces             & 378         & 339         & 270 \\
\quad won / tied / lost                & 28 / 9 / 63\% & 30 / 11 / 59\% & 77 / 23 / 0\% \\
attacks on hidden pieces, won / lost   & 43 / 43\%   & 44 / 42\%   & 39 / 47\% \\
\bottomrule
\end{tabular}
\end{center}
A network whose checkpoint failed to load (a fresh one, through the silent fallback of
\texttt{load\_model}) lost 66 to 67\% of its attacks on revealed pieces in the same games, so on
this decision both trained networks are close to an untrained one. In the rank-swap test (300
positions, the target shown as a revealed Scout or a revealed Marshal, everything else unchanged)
formal\_5's value of the attack rose by $+0.0016$ for the Scout, positive in 49\% of positions,
and formal\_4's by $+0.0011$, positive in 51\%, against a spread of the attack's value across
positions of 0.06 to 0.09. With a won exchange worth about three-quarters of an atom more than a
lost one in every return that contains it, 2750 episodes of material-only shaping did not make
the network read the revealed rank at all. A reward signal that clear and that frequent should
have been learnt, so the cause has to be sought in the learner or in what reaches it rather than
in the reward. None of the faults of the next sections has yet been shown to explain it.

\paragraph{Lesson.} When a change aimed at one failure leaves that failure exactly where it was,
measure the failure directly before reading the win rates: the probe settled in minutes what
three noisy evaluations could not. And a diagnostic that loads a checkpoint must fail loudly when
the load fails; here a silent fallback to a fresh network gave plausible-looking numbers.

% ------------------------------------------------------------------------------------------
\section{Asynchronous learner}

\paragraph{Change.} The gradient updates now run on a background thread with its own CUDA
stream while the main loop chooses moves and steps the parallel games. This is the actor-learner
split of Ape-X \cite{ApeX} and IMPALA \cite{IMPALA}, also used for large-scale asynchronous
reinforcement learning of language models \cite{AReaL}, here within a single process:
\begin{itemize}
  \item The replay ratio is unchanged: the loop still owes one update per
        \texttt{REPLAY\_UPDATE\_INTERVAL} learner moves and hands them to the learner. Once more
        than \texttt{ASYNC\_MAX\_LAG} $= 2$ updates are owed it waits (back-pressure), so the
        policy that acts is at most three updates behind the one being trained.
  \item Acting reads a copy of the network that is refreshed after every update. The learner
        must not share its network with the actor: the optimiser changes the weights in place,
        and NoisyNet acting switches the network's noise and train/eval mode.
  \item The replay memory, which both threads use, is behind a lock. Because the learner runs on
        its own CUDA stream, it waits for the actor's stream before sampling and finishes its
        reads before releasing the lock; otherwise a transition evicted by the actor could be
        freed by the allocator while the learner's kernels were still reading it.
  \item Anything that reads or replaces the learner's weights or optimiser state (checkpoints,
        copying the network into the second agent, the learning-rate schedule) pauses the
        learner first, and an error in the learner thread is re-raised in the training loop.
\end{itemize}

\paragraph{Result.} In two otherwise identical short runs (8 lanes, the cold start on, measured
from the first gradient update over about 10\,000 learner moves, on a GPU shared with another
training run) the synchronous trainer made 43.4 learner moves per second and the asynchronous one
59.3, i.e.\ 1.37 times as fast. Both performed exactly 1313 gradient updates, so the replay ratio
was the same.

\paragraph{Lesson.} Asynchrony is only safe if every piece of shared state is accounted for: the
weights (a separate acting copy), the replay memory (a lock), GPU memory on two streams (stream
ordering and synchronisation), and the replay ratio (back-pressure). The last one matters most
for reproducibility: without back-pressure an asynchronous learner changes the algorithm, not
just its speed.

% ------------------------------------------------------------------------------------------
\section{What the falling loss of formal\_4 did and did not show}

\paragraph{Observation.} Over the 3000 episodes of formal\_4 the C51 loss of Player~1 fell
steadily while its evaluations did not improve. Per block of 500 training episodes (about 11\,500
gradient updates each):
\begin{center}
\begin{tabular}{rcccc}
\toprule
Episodes & mean loss & mean $Q$ (action taken) & mean return & training games won (random / heuristic) \\
\midrule
0--500     & 3.63 & $-0.01$ & $-0.66$ & 34\% / 16\% \\
500--1000  & 3.42 & $-0.02$ & $-0.59$ & 39\% / 18\% \\
1000--1500 & 3.32 & $-0.06$ & $-0.63$ & 34\% / 18\% \\
1500--2000 & 3.21 & $-0.14$ & $-0.60$ & 36\% / 20\% \\
2000--2500 & 3.11 & $-0.22$ & $-0.57$ & 39\% / 21\% \\
2500--3000 & 3.05 & $-0.28$ & $-0.51$ & 43\% / 22\% \\
\bottomrule
\end{tabular}
\end{center}
The greedy evaluations over the same span won 40, 35, 30, 30, 38 and 43\% of the 40 games
against the random opponent and 25, 20, 10, 13, 15 and 8\% against the heuristic (episodes 500
to 3000). The run had the AAREN history planes switched off (\texttt{AAREN\_ENABLED} false: 30
input planes), so it says nothing about AAREN.

\paragraph{Reading.} The loss is the cross-entropy between the predicted return distribution and
a target built from the network's own target copy; a distribution spread evenly over the 51 atoms
scores $\ln 51 = 3.93$, about where the cold start left it (3.93). Most of the fall is the
network making its distributions sharper and moving its values towards what it actually gets:
the mean $Q$ of the actions taken went from $-0.01$ to $-0.28$, towards the mean return of about
$-0.55$, i.e.\ it learned more accurately that it was losing. Across blocks of 100 episodes the
loss and the training win rate correlated at $-0.49$, but the loss correlated with time at
$-0.97$; once a linear trend was removed from both, the correlation was $-0.14$. The training
games (with NoisyNet exploration on) did improve somewhat against the random opponent, from 34\%
to 43\% of about 215 games per block, roughly two standard errors; the greedy evaluations of 40
games did not show it.

\paragraph{Lesson.} A falling TD or distributional loss shows that the critic fits its own
moving targets better, not that the policy plays better; with a pessimistic, well-calibrated
critic it falls while the agent keeps losing. Two quantities that both drift with training time
correlate whatever their relation, so a loss--performance correlation should be computed on
detrended (or differenced) curves before it is read as coupling. Progress has to be judged on
evaluations against fixed opponents, with enough games for the difference to exceed its noise.

% ------------------------------------------------------------------------------------------
\section{Correcting the claims of the thesis chapters}

\paragraph{Change.} The Methodology and Results chapters were revised where the code or the data
did not support what they said. No reported number was changed; the readings of them were.
\begin{itemize}
  \item AAREN was described as trained end to end with the Q-network through a ``gradient
        bridge'', $\epsilon \cdot W_{proj}(\mathbf{0})$ added to the stored embeddings. An affine
        layer maps the zero vector to its bias, so the DQN loss reached only one of AAREN's 21
        parameter tensors, and none that depends on the move history. AAREN was trained, in
        effect, only by its supervised loss on revealed pieces, and the text now says so. In the
        thesis code its embedding was also taken at the padding step of every shorter history
        rather than at the history's last real step.
  \item The self-play win rate of 49.88\% (against 24.97\% for the DQN) was read as the agent
        doubling its strength. In self-play both sides are the same policy, and in both runs the
        decisive games split almost evenly (24.97/24.99\% and 49.88/43.43\%). The win rate doubled
        because the decisive share did (50\% to 93\%), and steps per win follows from that too.
  \item The loss--win correlation of $-0.787$ was read as the loss becoming a reliable proxy for
        play. It was computed on raw windowed curves. formal\_4 shows the raw correlation
        ($-0.49$) shrinking to $-0.14$ once a linear trend is removed from both curves, while its
        evaluations did not improve.
  \item AAREN's 24.23\% accuracy was compared only with a uniform 8.3\% baseline over 12 types.
        Scouts are 8 of the 40 pieces and 8 of the 33 movable ones, so always naming the most
        common type would score about 20 to 24\%.
  \item The logged average Q-values (0.70, 0.84, 1.84) were read as value resolution. They were
        the maximum over all 3600 actions, the artefact of the second entry of this appendix.
\end{itemize}
The validation table now reads ``Not tested'', ``Not supported'', ``Inconclusive'', ``Supported''
or ``Self-play only'' instead of ``Validated'' for every row. The Methodology chapter now lists all six
Rainbow components and separates the thesis configuration from that of the formal runs.

\paragraph{Why.} Re-reading formal\_4 against the chapters (previous entry). That run had AAREN
switched off, yet the chapters presented the falling loss and the self-play results as validation
of both the Rainbow learner and AAREN.

\paragraph{Lesson.} A metric from symmetric self-play measures the draw rate, not strength; a
correlation between two curves that both drift with time measures the drift; an accuracy needs
the majority-class baseline, not the uniform one; and a claimed gradient path should be checked
by listing which parameters actually receive a gradient. Each of these was cheap to check and
reversed a claimed validation.

% ------------------------------------------------------------------------------------------
\section{Replay priorities from the cross-entropy}

\paragraph{Observation.} The prioritized replay \cite{RainbowJMLR2021} sets a transition's priority to
$(\ell + 10^{-6})^{0.6}$, with $\ell$ the C51 loss of that transition, which is the cross-entropy
$-\sum_i m_i \log p_i$ between the projected target $m$ and the prediction $p$. The
cross-entropy is the KL divergence plus the entropy of the target,
$\mathrm{CE} = \mathrm{KL}(m \,\|\, p) + H(m)$. Rainbow prioritizes by the KL divergence, the part
that measures how wrong the prediction is. $H(m)$ is large for every non-terminal transition,
because its target is the network's own diffuse prediction one step on. It is at most $\ln 2$ for a
terminal one, whose target is the outcome itself, split over the two nearest atoms.

\paragraph{Measurement.} formal\_4's network of episode 3000 played 32 greedy games (against the
random opponent and the heuristic, 8 won, 16 lost and 8 drawn). They were recorded through the
trainer's 5-step and mirroring pipeline, which gave 13\,328 transitions; the rewards came from the
current material-only potential. For each transition:
\begin{center}
\begin{tabular}{lrccc}
\toprule
Transitions & count & CE & $H(m)$ & KL \\
\midrule
terminal                              & 320   & 4.15 & 0.51 & 3.64 \\
non-terminal, return with a battle    & 6790  & 3.76 & 3.73 & 0.02 \\
non-terminal, no battle               & 6218  & 3.42 & 3.40 & 0.02 \\
\bottomrule
\end{tabular}
\end{center}
With the cross-entropy, the 95th percentile of the priority was 1.2 times the 5th. The terminal
transitions, 2.4\% of the memory, held 2.6\% of the total priority, so they were sampled as
often as under uniform replay. With the KL divergence the ratio would be 10.3. The terminal
transitions would hold 40\% of the priority, about 16 times their share.

The terminal KL of 3.64 also shows that, one move before the end of a game, the network's
distribution put almost no mass near the actual outcome. Those 2.4\% of transitions are the only
ones whose target is the game result rather than the network's own estimate. Under the
cross-entropy priority the replay did nothing to make them count more.

A second, smaller fault: the 2.5-fold priority boost for battles is decided by the newest step of
a 5-step window. It therefore lands on the transition that \emph{starts} four moves before the
battle, not on the one that starts with the attack. The mirrored copy of a transition gets no
boost at all.

\paragraph{Change.} Made for formal\_6 (Section~\ref{sec:formal6}); at the time of this entry, none. The proposed fix is to compute the priority from
$\mathrm{CE} - H(m)$, and to decide the battle boost by the first step of the window. Whether this
changes play is not yet measured in a run.

\paragraph{Lesson.} A priority must measure surprise, not the size of the loss. With a
distributional loss the two differ by the entropy of the target, which here was almost the whole
loss and made prioritized replay uniform. The distribution of the priorities is cheap to check (the
ratio of high to low percentiles, and the share held by the transitions that carry the outcome),
and it should be checked once whenever the loss changes.

% ------------------------------------------------------------------------------------------
\section{An audit of the rest of the Rainbow learner}

\paragraph{Observation.} Before changing the replay priorities, the rest of the learner was read
and, where a doubt could be measured, probed on formal\_4's network of episode 3000.
\begin{itemize}
  \item \emph{Checked and correct:} the C51 projection, the Double DQN target over legal moves,
        the dueling combination, the factorised NoisyNet layers, the per-lane 5-step buffers
        (mirroring is applied after the $n$-step sum, so the two never mix), the flushing of
        open returns at the end of a game, the move-plus-reply transitions, the asynchronous
        learner's locking and weight copy, and the replay ratio (about 8 samples per stored
        transition, the usual DQN ratio).
  \item \emph{Instance normalisation and constant planes.} The backbone normalises every channel
        over the board, and planes 27--29 (observability, no-battle counter, game progress) are
        constant over the board, so they might have been normalised away. They are not. In 300
        fog-of-war positions, raising the no-battle counter from 0 to 0.99 lowered the network's
        value by 0.15 on average, as it should with a draw costing $-0.5$, and game progress
        by 0.11. Removing one own piece changed it by 0.02. A fresh network learnt a target
        from the counter plane equally well with and without the normalisation layers
        (correlation 0.99 in both cases).
  \item \emph{Dropout in the target.} The attention block has dropout 0.1 and the target network
        is never put in evaluation mode, so its targets carry dropout noise. Dropout alone moved
        the value of the greedy move by a standard deviation of 0.015, more than the NoisyNet
        noise (0.010), and about a third of the spread of the values across a position's legal
        moves (0.043).
  \item \emph{How much NoisyNet explores.} With the acting noise on, the move played differed
        from the greedy move in 4 to 6\% of positions (100 and 300 positions, 8 and 20 noise
        samples).
\end{itemize}

\paragraph{Change.} Made for formal\_6 (Section~\ref{sec:formal6}); at the time of this entry, none. The proposed changes are the evaluation mode for the target
network (or no dropout in the attention block) and the priority fix of the previous section.

\paragraph{Lesson.} A suspicion about an architecture is cheap to test: perturb the input it
concerns and see whether the output moves, or check whether a fresh network can fit a target
that depends only on it. Here the plausible suspicion (normalisation erasing the counters) was
wrong, and two quieter faults (an always-on dropout in the target, and priorities that do not
prioritise) were real.

% ------------------------------------------------------------------------------------------
\section{What AAREN can add to the observation}

\paragraph{Why it is off.} The formal runs set \texttt{AAREN\_ENABLED} to false, so the network sees
the 30 observation planes only. No document records the reason; the code enforces two
constraints. The cold start refuses AAREN, because a demonstration has no history of its own and
zero planes would teach the network what a game without history looks like. The asynchronous
learner also refuses it, because AAREN's supervised training shares the network. The plan was to
make the Rainbow learner work first and change one thing at a time.

\paragraph{Measurement.} The heuristic played 225 games as Player 1, alternately against the random
opponent and the heuristic. A history aggregator tracked Player 2's moves and labelled them at
battles exactly as the trainer does.
\begin{itemize}
  \item \emph{The ``has moved'' bit.} Of Player 1's attacks on a piece that had never moved,
        41\% (1472 attacks, against random) and 49\% (1028, against the heuristic) hit a Bomb
        or the Flag. Of its attacks on a piece that had moved, none did (2503 and 1712 attacks).
        The observation does not contain this bit, and the game state does not record which
        pieces have moved.
  \item \emph{Piece type from the history.} AAREN labels only pieces that moved at least twice.
        Its 6216 labelled reveals were split by game, 80\% for training and 20\% for testing,
        and a fresh AAREN was trained for 4000 batches. On the 1272 test reveals it named the
        type correctly in 35.6\% of cases. Always naming the most common type (Lieutenant) gave
        16.3\%, and ``a long move means a Scout, otherwise the most common type'' gave 23.9\%.
        Its accuracy was the same against the random opponent (34.9\%) as against the
        heuristic (36.0\%). Since the random opponent moves without regard to rank, most of
        what AAREN reads must be where a piece stood in the heuristic setup, rather than how it
        behaves. A setup prior is real information against this opponent pool, but it would
        shrink against varied setups.
\end{itemize}

\paragraph{Reading.} AAREN does carry information the observation lacks: it beats both simple rules
on held-out games. The largest single piece of history information, however, is whether a hidden
piece has moved at all. That bit separates a 45\% chance of attacking a Bomb or the Flag from
none, and it needs one observation plane rather than a 64-plane embedding. The thesis accuracy
of 24.23\% is close to the long-move rule here. Whether AAREN helps the \emph{agent} is a separate
question that only paired runs can answer, with and without it, on the same evaluations.

\paragraph{Change.} Made for formal\_6 (Section~\ref{sec:formal6}); at the time of this entry, none. Proposed order: add a ``has moved'' plane for enemy pieces (and mark
a Scout revealed by a long move), then compare AAREN on and off. The cold start and the
asynchronous learner would have to support AAREN for that comparison.

\paragraph{Lesson.} Before adding a learned module to supply missing information, find out what the
information is and how much of it a hand-made feature already captures. Here one bit carried most
of it, and a module trained on revealed pieces never sees the unmoved ones that bit is about.

% ------------------------------------------------------------------------------------------
\section{Moved-piece planes and three learner fixes (formal\_6)}
\label{sec:formal6}

\paragraph{Change.} Four changes, made together before formal\_6:
\begin{itemize}
  \item \emph{Two observation planes, 32 in all.} The environment now records, for every square,
        whether the piece on it has moved and whether it has ever moved more than one square at
        once, which only a Scout can do. Both facts are public: the opponent sees every move. The
        flags travel with the piece, including when it wins a battle, and are cleared when it is
        taken. Plane 30 marks the enemy pieces of unknown rank that have moved, so they are not a
        Bomb or the Flag. Plane 31 marks those that have made a long move, so they are Scouts.
        The GUI bot and the browser player build the same planes, and tests replay games in both
        and compare them plane by plane with the training environment.
  \item \emph{Replay priorities from the KL divergence.} A transition's priority is now
        $(\mathrm{KL}(m \,\|\, p) + 10^{-6})^{0.6}$, with
        $\mathrm{KL} = \mathrm{CE} - H(m)$, as in Rainbow \cite{RainbowJMLR2021}, instead of the
        cross-entropy. The loss itself is unchanged.
  \item \emph{The battle boost on the battle move.} The 2.5-fold priority of a battle now goes to
        the 5-step transition whose \emph{first} move was the attack. The mirrored copy of a
        transition now gets the same boosts as the original, for battles and for wins.
  \item \emph{No dropout in the target.} The target network is put in evaluation mode whenever
        it is copied, so its targets carry no dropout. Its NoisyNet noise stays on, as in
        Rainbow.
\end{itemize}
Everything else is formal\_5's configuration: material-only potential, cold start, asynchronous
learner, 16 lanes, 20\,000 episodes. The first convolution now reads 32 planes, so checkpoints of
formal\_4 and formal\_5 do not fit the new network; formal\_6 starts from zero like them.

\paragraph{Why.} The three preceding sections. The ``has moved'' bit separated a 41 to 49\% chance
that an attacked hidden piece was a Bomb or the Flag from none, and the observation did not contain
it. Under the cross-entropy priority the terminal transitions, the only ones whose target is the
game result, were sampled no more often than under uniform replay; under the KL divergence they
would hold about 16 times their share of the priority. The battle boost of the thesis code
landed four moves before the attack, and dropout moved the target value of the greedy move by
more than the NoisyNet noise did.

\paragraph{Result, cold start.} The cold start ran as before: the heuristic
won 50, lost 27 and drew 23 of its 100 demonstration games (39\,406 transitions), and 5000
pre-training updates took the loss from 4.13 to 3.95 and the margin loss to 0.076. The checks
set for formal\_6 were: whether attacks on
unmoved hidden pieces, and so losses to Bombs, become rarer than in formal\_4 and formal\_5;
whether the Q-value one move before the end of a game moves towards the outcome (formal\_4's
network put almost no mass there, a KL of 3.64); the spread of the priorities in the replay
memory, which should now be wide rather than flat; and whether the network finally reads the
rank of a revealed piece (the rank-swap test and the share of attacks on revealed pieces it
loses, 59 to 63\% for formal\_4 and formal\_5 against 0\% for the heuristic).

\paragraph{Result, to episode 9000.} The greedy evaluations of stage~1 (episodes 500 to 5000, 40
games each) won 17, 18, 10, 18, 11, 23, 14, 18, 15 and 13 games against the random opponent and
6, 5, 6, 13, 2, 7, 11, 8, 4 and 5 against the heuristic: no trend, and no better than formal\_4
or formal\_5. No gate was passed. At episode 5000 the stage's episode cap moved the opponent
schedule to stage~2, whose training opponents are the frozen and the smart heuristic only; the
random opponent and the plain heuristic drop out. The evaluations against the heuristic then won
7, 6, 7, 1, 5, 4, 1 and 7 of 40 (episodes 5500 to 9000). In the training games of stage~2 the
agent won about 17\% against the frozen heuristic and drew about 80\% against the smart one, and
games grew from about 350 to 550 moves, most of them ending in no-battle draws.

The probes of the previous runs were repeated on the networks of episodes 5000 and 9000, in 30
greedy games against the random opponent (and 300 positions for the rank-swap test):
\begin{center}
\begin{tabular}{lccc}
\toprule
 & ep.\ 5000 & ep.\ 9000 & heuristic \\
\midrule
games won / lost / drawn                      & 11 / 15 / 4 & 0 / 3 / 27 & 23 / 4 / 3 \\
attacks per game                              & 34          & 9          & 39 \\
attacks on revealed pieces lost               & 66\%        & 73\%       & 0\% \\
rank-swap: Scout $-$ Marshal, positive in     & 48\%        & 33\%       & -- \\
share of attacks on unmoved hidden pieces     & 27\%        & 0\%        & 31\% \\
\quad of which Bomb or Flag                   & 44\%        & --         & 40\% \\
\bottomrule
\end{tabular}
\end{center}
formal\_5's network (episode 2750, 30 planes) put 30\% of its attacks on unmoved hidden pieces,
37\% of them Bombs or the Flag. So by episode 5000 the moved-piece planes had not changed whom the
agent attacks, and it still did not read the revealed rank (the Scout and Marshal versions of an
attack differed by $-0.001$ on average). By episode 9000 it no longer attacked unmoved pieces at
all, but it attacked hardly anything: it drew 27 of 30 games against the random opponent, which
formal\_6 had beaten 11 times at episode 5000. Its value of an attack was about $-0.45$ whatever the
target, close to the $-0.5$ of a draw.

The C51 loss reported by the trainer fell from 2.47 (first 500 episodes) to 0.38 (last 500), far
below the 3.0 to 3.7 of formal\_4 and formal\_5. The reported loss is the cross-entropy weighted by
the importance-sampling weights of the prioritized replay, and those weights change with the
priorities, so the number is not comparable across the change of priority and shows nothing about
play on its own.

\paragraph{Reading.} Two separate things. In stage~1, the moved-piece planes and the learner fixes
changed nothing that the probes measure: the network still does not learn the outcome of a battle
from what it can see, the failure of formal\_4 and formal\_5. In stage~2, a change of opponents
alone, with the environment unchanged, turned an agent that beat the random opponent about a third
of the time into one that only draws: the agent trained against a pool in which it mostly draws
(the smart heuristic) or loses (the frozen one), its values moved to the value of a draw, and it
stopped attacking. This is the forgetting that the removal of the curriculum's environment changes
was meant to avoid, caused here by the opponent schedule that was kept.

\paragraph{Lesson.} An episode cap that moves an agent to the next stage of an opponent schedule
when it has not passed the gate turns the schedule into a fixed timetable, and dropping the weak
opponents from the pool removes the games the agent can win; with only draws and losses left, a
value learner learns that nothing it does matters and plays for the draw. Keep the earlier
opponents in the pool, and evaluate against them throughout, so that forgetting shows up when it
starts. (formal\_7 kept them, and the collapse happened anyway, so dropping them was not the cause:
Section~\ref{sec:formal7}.)

\paragraph{Lesson.} When a run is stopped to fix several faults at once, the next run cannot say
which fix mattered; it can only say whether they did together. Here that trade was accepted
because each fix restores the standard algorithm rather than adding something new, and the
paired comparison that matters more (AAREN on and off) still has to be run with one difference.

% ------------------------------------------------------------------------------------------
\section{Why the network does not learn the rank of a revealed piece}
\label{sec:revealed-signal}

\paragraph{Observation.} In every formal run the agent lost most of its attacks on enemy pieces
whose rank it could see (59 to 73\%), and its value of an attack did not depend on that rank
(Section~\ref{sec:formal6}). Three explanations were possible: the information does not reach
the training targets, the network cannot represent the rule, or the training signal is too weak
for the learner to pick up. Three probes on formal\_6's network of episode 5000 separate them.

\paragraph{Measurement 1: the targets.} The network played 32 greedy games (6 won, 22 lost, 4
drawn) through the trainer's pipeline, which stored 11\,426 transitions. Each one was classed by
its first move; an attack on a revealed piece was classed won, tied or lost by the battle rule.
For the non-terminal ones:
\begin{center}
\small
\begin{tabular}{lrrccc}
\toprule
First move & count & share & 5-step reward & $\mathbb{E}[m]$ & network $Q$ \\
\midrule
quiet move                   & 9524 & 83.4\% & $-0.016$ & $-0.110$ & $-0.108$ \\
attack on a hidden piece     & 1304 & 11.4\% & $-0.002$ & $-0.035$ & $-0.017$ \\
attack on a revealed piece, won  & 146 & 1.3\% & $+0.023$ & $-0.040$ & $-0.079$ \\
attack on a revealed piece, lost & 388 & 3.4\% & $-0.027$ & $-0.104$ & $-0.074$ \\
\bottomrule
\end{tabular}
\end{center}
The targets do carry the outcome of the exchange: their means differ by 0.064 between won and
lost attacks. The network's $Q$ differs by $-0.005$. The agent also chose attacks on revealed pieces
that it lost 2.7 times as often as ones it won. Each target distribution $m$ is wide, with a standard
deviation of 0.77 for these transitions, because it carries the uncertainty of an outcome hundreds
of moves away. The target \emph{means} of one class varied much less across transitions (0.13).

\paragraph{Measurement 2: capacity.} A fresh network of the same size was trained on 3267
positions from these games. In each, one attack on an adjacent enemy piece was chosen, and the
target square was edited to a revealed piece of random rank. The C51 target was a function of
the battle outcome $o \in \{1, 0, -1\}$ of that attack alone. Every batch of 128 consisted of such
attacks. The measure is the gap between the network's values of the attack when the target is
shown as a Scout (a win) and as a Marshal (a loss), over 300 held-out positions:
\begin{center}
\small
\begin{tabular}{p{7.2cm}lcc}
\toprule
Target & l.r. & Scout $>$ Marshal & gap (ideal) \\
\midrule
$0.5\,o$, no noise & $3\cdot 10^{-5}$ & 100\% & 0.44 (1.00) \\
$0.5\,o$, no noise & $10^{-3}$ & 100\% & 0.97 (1.00) \\
$0.07\,o$ plus sampled noise of s.d.\ 0.77 & $3\cdot 10^{-5}$ & 68\% & 0.003 (0.14) \\
\quad the same with signal $0.35\,o$ & $3\cdot 10^{-5}$ & 94\% & 0.023 (0.70) \\
\quad the same with 12 ``attack wins / ties / loses'' planes & $3\cdot 10^{-5}$ & 72\% & 0.005 (0.14) \\
distributions of s.d.\ 0.77, mean $0.07\,o$ plus noise of s.d.\ 0.13 & $3\cdot 10^{-5}$ & 99\% & 0.028 (0.14) \\
\quad the same plus a squared error on the mean & $3\cdot 10^{-5}$ & 99\% & 0.052 (0.14) \\
\bottomrule
\end{tabular}
\end{center}
(After 2000 batches. ``Scout $>$ Marshal'' is the share of held-out positions in which the attack
is valued higher with the target shown as a Scout.) With a clean target the network learnt the
rule within 1000 batches, even at the run's learning
rate, so it can read a revealed rank and compare it with the attacker's. Planes that state the
outcome of every possible attack outright hardly helped, so finding the relation is not the
difficulty either. What decides is the size of the signal relative to the width of the target. The
win--loss difference of the material potential is under one atom of the support (0.08 of $\pm 2$
over 51 atoms) and a twelfth of the target's width. Under the cross-entropy, a shift of the mean
that small changes the loss by little, so it is learnt slowly. With targets shaped like the real
ones, 2000 batches of nothing but these attacks learnt a fifth of the gap. Adding a squared error
between the predicted and target means doubled that.

\paragraph{Reading.} The rank information reaches the network, and the network can use it. The
training signal for using it is small and diluted. Attacks on revealed pieces were 5\% of the stored
transitions, so the real learner sees about 6 in a batch of 128. Those 6 compete with everything
else the network fits, against a target network that changes every 1000 updates. The probes show
the signal is learnable when concentrated. That dilution and the moving targets account for the
rest is an inference, not a measurement. This is also why the earlier fixes did not change the
behaviour. The material-only potential made the signal clean, but no larger than a fraction of an
atom. The moved-piece planes and the priority fixes do not touch it. Nor would AAREN: it supplies
information about hidden pieces, while this failure concerns pieces that are already revealed.

\paragraph{Change.} The first and third were made for formal\_7 (Section~\ref{sec:formal7}). The
options the probes support were, by cost: an expected-value loss beside the C51 loss (it doubled
the learnt gap); a larger weight on the material potential with a wider value support (a five-fold
signal got the ordering right in 94\% of positions); and a priority boost for attacks on revealed
pieces, which raises their share of the batch.

\paragraph{Lesson.} A distributional critic learns differences in the mean slowly when they are
small next to the width of the return distribution, even when they are exact. Before relying on
a reward term to teach a decision, compare that term's difference between the right and the wrong
choice with the spread of the targets. Then check, on a fresh network trained on that decision
alone, how much of the difference it learns in a realistic number of updates.

% ------------------------------------------------------------------------------------------
\section{A mean loss, revealed attacks in the batch, and the opponent pool (formal\_7)}
\label{sec:formal7}

\paragraph{Change.} formal\_6 ended at episode 11\,680, when its process was killed from outside
(its log shows no error; it was started from the desktop app's shell, and formal\_7 is started
detached from it). Its evaluations against the heuristic from episode 9500 on were 5, 6, 8, 5 and
5 wins of 40, with no gate passed. formal\_7
starts from zero with formal\_6's configuration and three changes:
\begin{itemize}
  \item \emph{A loss on the mean.} The squared error between the means of the predicted and the
        target distribution, $\bigl(\sum_i p_i z_i - \sum_i m_i z_i\bigr)^2$ with weight 1, is
        added to each transition's C51 cross-entropy. Both are weighted by the importance-sampling
        weight. The priorities are still the KL divergence. The loss reported by the trainer now
        includes this term, so it is not comparable with that of earlier runs.
  \item \emph{Attacks on revealed pieces in the batch.} A transition whose first move attacks an
        enemy piece of known rank, and its mirror image, keeps four times its priority through
        every priority update. The importance-sampling weights leave this factor uncorrected, so
        these transitions count four times as much in the loss, not just four times as often.
        At about 5\% of the replay, they become about 17\% of each batch.
  \item \emph{The opponent pool after stage 1.} Stages 2 and 3 keep the random opponent and the
        heuristic in their pools beside the frozen and smart heuristics (or the self-play
        snapshot). The usual weighting (more games against the opponents it beats less often)
        still gives the random opponent at least 15\%. Every evaluation now also plays the random
        opponent, so forgetting shows as soon as it starts. Stage 1 and the episode caps are
        unchanged.
\end{itemize}

\paragraph{Why.} Section~\ref{sec:revealed-signal}. The win--loss difference of an exchange is
present in the targets but is a twelfth of their width, so the cross-entropy learns it slowly.
A mean loss doubled the learned gap in the probe. The attacks that carry the difference were 5\%
of each batch. Section~\ref{sec:formal6}: once the stage-1 cap moved formal\_6 to a pool it could
only draw or lose against, it stopped attacking.

\paragraph{Result, to episode 2500.} The cold start ran as in formal\_6 (the heuristic won 53,
lost 27 and drew 20 of the 100 games; 5000 pre-training updates, loss 4.17 to 3.94 with the mean
term, margin loss 0.074). Wins in the 40-game evaluations, with formal\_6 at the same episodes in
brackets:
\begin{center}
\begin{tabular}{lccccc}
\toprule
Episode & 500 & 1000 & 1500 & 2000 & 2500 \\
\midrule
random    & 16 (17) & 13 (18) & 12 (10) & 23 (18) & 17 (11) \\
heuristic & 4 (6)   & 4 (5)   & 10 (6)  & 7 (13)  & 7 (2) \\
\bottomrule
\end{tabular}
\end{center}
Summed over the five evaluations, formal\_7 won 81 games against the random opponent to formal\_6's
74, and 32 against the heuristic to 32. One evaluation of 40 games has a standard error of about
3 wins at these rates, so stage 1 shows no difference yet. That is expected: two of the three
changes act on stage 2 and later, and the mean loss acts on a signal the evaluations measure only
indirectly. The test is still the rank-swap probe and the share of revealed attacks lost, and,
after episode 5000, whether the evaluations against the random opponent hold up.

\paragraph{Result, the revealed-rank probe at episodes 2500 and 3500.} Each checkpoint of both
runs played 32 greedy games (same seed). The attacks on revealed pieces in them were classed by
the battle rule, and the rank swap of Section~\ref{sec:revealed-signal} was run on one fixed set
of 300 positions for all four checkpoints:
\begin{center}
\small
\begin{tabular}{lcccc}
\toprule
 & formal\_6, 2500 & formal\_7, 2500 & formal\_6, 3500 & formal\_7, 3500 \\
\midrule
attacks on revealed pieces & 478 & 484 & 512 & 470 \\
\quad won / tied / lost & 26 / 10 / 64\% & 32 / 6 / 62\% & 30 / 9 / 62\% & 29 / 8 / 63\% \\
Scout $>$ Marshal & 65\% & 64\% & 66\% & 67\% \\
mean gap & 0.021 & 0.020 & 0.032 & 0.022 \\
$\mathbb{E}[m]$ gap, won $-$ lost & 0.088 & 0.111 & 0.101 & 0.115 \\
\bottomrule
\end{tabular}
\end{center}
By episode 3500 (about 650\,000 updates) the two runs do not differ. Both still lose about
two-thirds of the attacks they choose on pieces whose rank they can see, and both rank the Scout
above the Marshal in two-thirds of the positions, with a gap of a fifth of the target's 0.1. The
targets carry the difference as before. So the mean loss and the four-fold priority have not yet
made the network learn it faster in the real run. The probe had predicted twice the learning
speed on a batch made only of these attacks. In the run they are about a sixth of the batch,
and the targets move. The value of every attack was also lower in formal\_7 at episode 3500
($-0.22$ to $-0.26$, against $-0.05$ to $-0.12$ in formal\_6). That checkpoint comes after the
resume, which emptied the replay of its own games, so the cause is not clear.

\paragraph{Result, to episode 5500: stage 2 collapses as before.} Stage 1 ended at its cap of
5000 episodes. Its last evaluations were 15, 18 and 17 wins of 40 against the random opponent
and 7, 6 and 6 against the heuristic (episodes 4000 to 5000), as in formal\_6. 500 episodes into
stage 2, the greedy agent drew 37 of 40 games against the random opponent, won none and lost 3.
It won 7 of 40 against the heuristic. Episodes slowed from about 1.9 to 6.9 seconds, as the
games ran on to the draw rules. The random opponent was in the stage-2 pool: it got 112 of the
500 training games. The agent won the first of them about half the time and then stopped winning
against it. So keeping the opponents it can beat in the pool did not prevent the collapse;
Section~\ref{sec:formal6}'s explanation, that the agent forgot them, is wrong.

Only the opponents change at the stage switch: the game, its limits and the shaping stay the
same. A reading that fits the numbers is that stage 2 made avoiding battle the better policy. Against
the frozen and smart heuristics, three-quarters of the stage-2 pool, the agent won about 15\% and
lost about 70\% of its training games. A game it fought out was then worth about
$0.15 - 0.70 = -0.55$, no more than a draw ($-0.5$). An exchange it lost also cost material
shaping at once, while avoiding battle cost nothing until the game ended. The observation does
not say which opponent the agent is playing, so the same passive policy was played against the
random opponent. The underlying weakness is the one measured above: it loses two-thirds of the
attacks it chooses on pieces whose rank it can see. Until that changes, fighting does not pay
against a competent opponent.

\paragraph{Lesson.} When a curriculum moves the agent to opponents it mostly loses to, check the
value of a draw against the value of playing on. If the draw is worth as much, the agent will learn
to draw, and it will do so against every opponent it cannot tell apart from them.

The process was killed at about episode 3000 by a console window-close event
(Section~\ref{sec:engineering}). It was resumed from the checkpoint of episode 2750 with the
networks, the target network, the optimiser and the curriculum state, but not the replay
buffer. On resuming, the trainer played 100 new heuristic games (48 won, 24 lost, 28 drawn;
36\,214 demonstration transitions) to refill it, and skipped the pre-training updates. Its own
transitions of episodes 1 to 2750 are gone, so for a while after the resume the replay holds a
larger share of demonstrations than it did before. Results after episode 2750 are of a resumed
run.

\paragraph{Lesson.} When a probe shows that a learner cannot pick up a small but exact signal, the
fix should act on that signal, not on the agent's inputs. Two levers do: the loss, so the signal
weighs more against the noise it competes with, and the data, so the cases that carry it are a
larger share of each batch.

% ------------------------------------------------------------------------------------------
\section{What the agent has to learn, and in what order}
\label{sec:skill-ladder}

\paragraph{Observation.} Read together, the formal runs show the agent stuck at one step of a
sequence of skills that Stratego asks for. Each skill pays off only through the one before it,
and the opponent schedule had assumed skills the agent did not have.
\begin{enumerate}
  \item \emph{The rules.} Which moves are legal. The action mask supplies this; nothing has
        to be learned.
  \item \emph{Engagement.} Moving pieces forward and making contact, instead of shuffling until
        a draw rule ends the game. Learned in every run since the draw was given a cost: in stage 1
        the agent attacks, and it wins 30 to 58\% of its evaluation games against the random
        opponent, about 40\% on average.
  \item \emph{Comparing two known ranks.} Attacking a revealed piece it beats and leaving or
        fleeing one that beats it. This is the whole game under full information, and it needs
        no inference. It is \emph{not} learned: from formal\_4 to formal\_7 the agent lost 59 to
        73\% of the attacks it chose on revealed pieces, and its values barely depend on the rank it
        sees (Sections~\ref{sec:revealed-signal} and~\ref{sec:formal7}). An agent with this skill
        should beat a random opponent nearly every time, because the random opponent offers its
        pieces at random. Winning 40\% of those games is itself a sign that the skill is missing.
  \item \emph{Inference about hidden pieces.} A piece that has moved is neither Bomb nor Flag. A
        piece that moved more than one square is a Scout. A low piece can be sent to make a
        strong one reveal itself. formal\_6 gave the network these facts as planes 30 and 31.
        This skill pays off only through skill~3: knowing more about a piece is worth something
        only to an agent that attacks the right pieces once it knows.
  \item \emph{Managing information.} Keeping its own strong pieces hidden, protecting revealed
        ones, and the Spy against the Marshal. It pays off by denying the opponent skills 3
        and 4.
  \item \emph{Strategy.} Finding and defending the Flag, tempo, the setup and the endgame. Only
        the outcome rewards it.
\end{enumerate}

\paragraph{Reading.} The reward signal for each skill weakens down the list. The material
potential rewards skill~3 at once, as each exchange is won or lost. Skill~4 gets no immediate
reward; it is paid later, through better exchanges. Skill~6 is paid only by the outcome, hundreds
of moves later. This is why skill~3 is the bottleneck. It has the most direct signal of the skills
that have to be learned, and everything above it is built on it. (Why that signal is not learned
is still open. It is not only its size: formal\_8 made it two and a half times larger in the
targets, and the network still did not respond to the rank; Section~\ref{sec:formal8}.)

The opponent schedule assumed the skills in a different order. Stage 2's frozen and smart
heuristics play skills 3 to 5 well, so playing them presumes the agent can already compare ranks.
The stage-1 gates measured win rates against random and the heuristic, which mix every skill
together. And the episode cap moved the agent on whether or not it had passed them. An agent
without skill~3 that meets opponents with it loses most exchanges. It can then do no better than
draw, which is what formal\_6 and formal\_7 learned to do.

\paragraph{Lesson.} Name the skills a game requires, the order in which their rewards depend on
each other, and a direct measure of each. Then advance a curriculum on the measure of the skill the
next stage needs, not on a win rate that mixes them all. The direct measure of skill~3 is the
share of attacks on revealed pieces that the agent loses in its own games. The evaluations report
it from formal\_8 on.

% ------------------------------------------------------------------------------------------
\section{A five-fold material signal and stages without a cap (formal\_8)}
\label{sec:formal8}

\paragraph{Change.} formal\_7 was stopped at episode 5690, in stage 2. formal\_8 starts from zero
with formal\_7's configuration (mean loss, priority boost on attacks on revealed pieces, earlier
opponents kept in the pool) and three changes:
\begin{itemize}
  \item \emph{The material potential is five times larger} (weight 5 instead of 1). The support of
        the value distribution is widened from $\pm 2$ to $\pm 6$ (0.24 per atom) so that it still
        covers every return (outcome plus $-\Phi(s_0)$ with $\Phi \in [-5, 5]$). A won exchange
        of average size now moves the target by about 0.3, more than an atom, against the same
        outcome noise as before.
  \item \emph{No episode cap on any stage.} A stage ends only when its evaluation gates are
        passed (stage 1: 70\% against the random opponent and 50\% against the heuristic).
  \item \emph{A direct measure of skill~3 in the evaluations.} Every evaluation game counts the
        agent's attacks on revealed pieces and how many it won or lost by the battle rule; the log
        shows the share lost beside each opponent's result.
\end{itemize}

\paragraph{Why.} Section~\ref{sec:skill-ladder}. The capacity probe of
Section~\ref{sec:revealed-signal} put the rank-swap ordering right in 94\% of positions when the
signal was five times larger, against 68\% at the run's size. Stage 1 is the only stage in which
no run has collapsed, so the agent stays there until it has the skill the next stage requires.

\paragraph{Caveat.} The shaping is potential-based, so it leaves the optimal policy unchanged
\cite{NgPBRS}. The optimal shaped values are the unshaped ones minus $\Phi(s)$, a term that
depends only on the state, so in the limit the gap between two actions is the same at any weight.
A larger potential makes the exchange visible in the early targets, before the bootstrapped values
have learned $-\Phi$. It speeds up learning and does not change what is learned. This is the
intended effect: the probes showed the agent learning skill~3 too slowly, not learning the wrong
thing.

\paragraph{Result.} The cold start ran as before (the heuristic won 51, lost 23 and drew 26 of
the 100 games; 5000 pre-training updates, loss 4.40 to 4.00, margin loss 0.084). The run was
stopped at episode 2386. Its evaluations, with formal\_7 at the same episodes in brackets:
\begin{center}
\small
\begin{tabular}{lccccc}
\toprule
Episode & 500 & 1000 & 1500 & 2000 & sum \\
\midrule
wins against random    & 10 (16) & 11 (13) & 18 (12) & 18 (23) & 57 (64) \\
wins against heuristic & 4 (4)   & 12 (4)  & 5 (10)  & 5 (7)   & 26 (25) \\
revealed attacks lost, against random    & 63\% & 69\% & 62\% & 60\% & \\
revealed attacks lost, against heuristic & 77\% & 72\% & 72\% & 78\% & \\
\bottomrule
\end{tabular}
\end{center}
The wins do not differ from formal\_7's. The new measure of skill~3 does not improve over the
run. In each 40-game evaluation the agent attacked 450 to 510 revealed pieces against the random
opponent and 160 to 180 against the heuristic, and it lost 60 to 78\% of those attacks.

The probe of Section~\ref{sec:formal7} on the checkpoint of episode 2000 separates the target from
what the network learned. The won--lost gap of the target means was 0.25 (about 0.1 in formal\_7):
the five-fold potential did make the signal in the targets two and a half times larger. The
network's $Q$ gap between won and lost attacks was 0.06. In the rank swap the network valued
the attack higher on a revealed Scout than on a revealed Marshal in 53\% of the 300 positions,
with a mean gap of 0.003. That is no better than chance, and below formal\_7's 66\% at the same
episode. (The probe ran with formal\_8's support and potential. Only the 66\% is comparable for the
formal\_7 checkpoint; its values are on the wrong scale.)

So the size of the signal relative to the target's width was not what kept the network from
learning skill~3. Section~\ref{sec:revealed-signal} explained the failure that way, from a fresh
network trained on batches made only of these attacks; that network learned the rule. In the
run, a signal two and a half times larger left the network's response to the rank at chance.
Something that differs between the probe and the run stops it: the mix of the batch, the moving
bootstrapped targets, the NoisyNet head, or the rest of the training pipeline. Which one is the
next question. A direct test is to train the probe's fresh network on batches drawn from the
run's own replay buffer, adding one difference at a time, and watch the rank swap.

\paragraph{Lesson.} A probe that isolates one factor can show that the factor is \emph{sufficient}
to explain a failure, not that it is the explanation. The capacity probe showed a small signal
learnt slowly; the run with a larger signal learnt nothing more. Before scaling a fix up to a full
run, reproduce the failure inside the probe by adding back the run's conditions one at a time,
until the probe fails as the run does.

% ------------------------------------------------------------------------------------------
\section{Engineering notes}
\label{sec:engineering}

\begin{itemize}
  \item formal\_6 and formal\_7 were both killed partway, with no Python error. formal\_7's error
        log showed \texttt{forrtl: error (200): program aborting due to window-CLOSE event} once
        for each process. The Intel Fortran runtime that numpy's maths library loads installs a
        console handler that aborts the process when its console window closes. Starting the run
        detached from the app's shell was not enough, because the launcher still gave it a
        console. Runs are now started with \texttt{pythonw.exe}, which has no console, with the
        output sent to log files and \texttt{FOR\_DISABLE\_CONSOLE\_CTRL\_HANDLER=1} set. A
        long run should not own a console window that something else can close.
  \item An \texttt{import time} statement inside the training loop made \texttt{time} a local
        name for the whole training function, so any use of \texttt{time} before that line
        failed. Modules should be imported once, at the top of the file.
  \item A checkpoint loader that silently falls back to a fresh network turns a missing file into
        a misleading result: a probe of a checkpoint that did not exist yet reported the values of
        an untrained network. Diagnostic scripts now assert that the file exists.
  \item Most source files have Windows line endings. An edit script that matches multi-line text
        must normalise them first, or it silently matches nothing; a patch helper that asserts the
        number of matches made every such miss fail loudly instead.
  \item The first run of the capacity probe above passed an all-zero next-state distribution to
        the C51 projection for its terminal targets. The projection moves probability mass, so
        the targets were empty, the loss was zero, and no network learnt anything. Every
        configuration then gave the same numbers to three digits, which exposed it. A probe
        should assert that its targets are distributions (they sum to one) before trusting a
        negative result.
  \item A patch helper made tolerant of trailing whitespace also treated the end of a replaced
        block as optional whitespace, and so ate the indentation of the line that followed. Four
        source files no longer compiled. Compiling every edited file at once, before running any
        test, found all of them in one pass.
\end{itemize}
